\def\lettersWithEntries{ABCDEFGHIJLMNOPQRSTV}
% ==================================================
%----------------------------------------------------------------------------------------
%	SECTION A
%----------------------------------------------------------------------------------------

\section*{A}

\begin{multicols}{3}

\extramarks{A}{A}
\entry{À dessein}{}{}{Intentionnellement.}{}{}

\entry{À l’emporte-pièce}{}{}{De façon brutale, tranchée ou sans nuance.}{}{}

\entry{À la fleur de l’âge}{}{}{À l’époque de la jeunesse, dans la période de pleine vigueur.}{}{Un bon petit diable à la fleur de l'âge, La jambe légère et l'œil polisson}

\entry{À la hussarde}{}{}{Brutalement, sans précaution ni délicatesse.}{}{}

\entry{À la va-comme-je-te-pousse}{}{}{Sans méthode, de façon improvisée et désordonnée.}{}{}

\entry{À tout berzingue}{}{}{Très vite.}{}{}

\entry{À vau-l’eau}{}{}{En se dégradant, sans direction ni contrôle.}{}{}

\entry{Affliction}{}{}{À vérifier dans un dictionnaire lexicographique.}{}{Disons que dans son affliction, Il trouve des compensations.}

\entry{Aïeux}{}{}{À vérifier dans un dictionnaire lexicographique.}{}{Je serais honteux de mon sang, Des aïeux de qui je descends}

\entry{Alangui}{}{}{À vérifier dans un dictionnaire lexicographique.}{}{Peu de chances qu'on voie mes belles odalisques, Déposer en grand deuil au pied de l'obélisque, Quelques gerbes de fleurs. Tout au plus gentiment diront-elles : « Peuchère, Le vieux Priape est mort », et, la cuisse légère, Le regard alangui,}

\entry{Ambages}{}{}{À vérifier dans un dictionnaire lexicographique.}{}{Sans ambages, Dans les pages, Du missel,}

\entry{Argousins}{}{}{À vérifier dans un dictionnaire lexicographique.}{}{Qu'il avait un petit cousin, au gué, au gué, Haut placé chez les argousins, au gué, au gué}

\entry{Au diable Vauvert}{}{}{Très loin, dans un endroit éloigné.}{}{Emportent les trépassés jusqu'au diable vauvert}

\entry{Aubade}{}{nom}{Musique ou chant donné au petit matin, traditionnellement sous les fenêtres de quelqu’un.}{}{J'ai graissé la patte au berger, Pour lui faire jouer une aubade.}

\entry{Avecques}{}{}{À vérifier dans un dictionnaire lexicographique.}{}{D'avecques ma femme, J'ai foutu le camp}

\entry{Avoir l’âme en peine}{}{}{Être triste, tourmenté ou malheureux.}{}{}

\entry{Avoir le cœur en fête}{}{}{Être très heureux.}{Être très heureux.}{}

\entry{Avoir les crocs}{}{}{Avoir très faim.}{}{}

\entry{Avoir les yeux de Chimène}{}{}{Regarder quelqu’un avec amour ou admiration.}{}{}

\entry{Avoir maille à partir}{}{}{Avoir un différend, une querelle avec quelqu’un.}{}{}

\entry{Avoir pignon sur rue}{}{}{Être bien établi, connu et respectable.}{}{}

\entry{Avoir une âme en peine}{}{}{Être profondément triste ou affligé.}{}{}

\end{multicols}

%----------------------------------------------------------------------------------------
%	SECTION B
%----------------------------------------------------------------------------------------

\section*{B}

\begin{multicols}{3}

\extramarks{B}{B}
\entry{Battre sa coulpe}{}{}{Reconnaître sa faute, s’en accuser.}{}{}

\entry{Bedeau}{}{nom}{Laïc chargé de certaines fonctions matérielles dans une église.}{}{L'maître d'école et ses potaches, le maire, le bedeau, le bougnat, négligeaient carrément leur tâche pour voir ça}

\entry{Bélître}{}{}{À vérifier dans un dictionnaire lexicographique.}{}{On murmure in petto : « C'est un vrai Nicodème, Un balourd, un bélître, un bel âne bâté. »}

\entry{Béotien}{}{nom/adj.}{Personne ignorante ou peu cultivée dans un domaine; à l’origine, habitant de Béotie.}{}{Hélas un béotien à la place du sien, lui proposa son blase}

\entry{Bibine}{}{nom}{Bière, généralement bon marché, dans un emploi populaire.}{}{L'éponge trempée dans Dieu sait quelle bibine,}

\entry{Bigot}{}{nom/adj.}{Personne attachée de manière étroite et excessive aux pratiques religieuses.}{}{J'ai perdu la tramontane, En perdant Margot, Qui épousa, contre son âme, Un triste bigot...}

\entry{Billevesées}{}{}{À vérifier dans un dictionnaire lexicographique.}{}{Et les billevesées de tous les répertoires, Et les morts pour que naisse un avenir plus beau,}

\entry{Bon gré mal gré}{}{}{Que cela plaise ou non.}{}{}

\entry{Bordel}{}{}{À vérifier dans un dictionnaire lexicographique.}{}{Sitôt privée de ma tutelle, Ma pauvre amie, Courut essuyer du bordel, Les infamies}

\entry{Bougnat}{}{nom}{Marchand de charbon et, traditionnellement, débitant de boissons d’origine auvergnate.}{}{L'maître d'école et ses potaches, le maire, le bedeau, le bougnat, négligeaient carrément leur tâche pour voir ça}

\entry{Bourrelé}{}{adj.}{Qui est tourmenté ou accablé par un sentiment, notamment le remords.}{}{Le criminel saisi d'un zèle expiatoire, Qui bat sa coulpe bourrelé par le remords,}

\entry{Braire}{}{verbe}{Crier, en parlant de l’âne; par extension, crier d’une manière désagréable.}{}{Elle eut beau pousser des sanglots, Braire à tue-tête, Comme je n'étais qu'un salaud, J'me fis honnête}

\entry{Branque}{}{}{À vérifier dans un dictionnaire lexicographique.}{}{Qu'ils aient comme ce branque compté la musique pour moins que zéro}

\entry{Brummel}{}{}{À vérifier dans un dictionnaire lexicographique.}{}{Tous les Brummel, les dandys, les gandins, il les considérait avec dédain}

\entry{Bucolique}{}{}{À vérifier dans un dictionnaire lexicographique.}{}{Une adorable bucolique, Il est mélancolique}

\end{multicols}

%----------------------------------------------------------------------------------------
%	SECTION C
%----------------------------------------------------------------------------------------

\section*{C}

\begin{multicols}{3}

\extramarks{C}{C}
\entry{Callipyge}{}{adj.}{Qui a de belles fesses, bien proportionnées.}{}{VENUS CALLIPYGE}

\entry{Calotins}{}{}{À vérifier dans un dictionnaire lexicographique.}{}{Ils ne savent pas ce qu'ils perdent, Tous ces fichus calotins, Sans le latin, sans le latin, La messe nous emmerde}

\entry{Camarde}{}{}{À vérifier dans un dictionnaire lexicographique.}{}{Et Dieu fit signe à la camarde, De l'expédier rue Froidevaux...}

\entry{Chaumine}{}{}{À vérifier dans un dictionnaire lexicographique.}{}{Vers son clocher, sa chaumine, Ses parents et sa promise}

\entry{Chenue}{}{}{À vérifier dans un dictionnaire lexicographique.}{}{Il avait la tête chenue, Et le cœur ingénu}

\entry{Chevet}{}{}{À vérifier dans un dictionnaire lexicographique.}{}{Ne vous étonnez pas, ma chère, Si vous trouvez, Les vers de jadis et naguère, A mon chevet.}

\entry{Coi}{}{adj.}{Silencieux, sans parler; immobile et muet.}{}{Qu'je m'démène ou qu'je reste coi, Je passe pour un je-ne-sais-quoi !}

\entry{Complainte}{}{}{À vérifier dans un dictionnaire lexicographique.}{}{Dieu fasse que ma complainte aille, tambour battant, Lui parler de la pluie, lui parler du gros temps}

\entry{Cornegidouilles}{}{}{À vérifier dans un dictionnaire lexicographique.}{}{Tous les morbleus, tous les ventrebleus, Les sacrebleus et les cornegidouilles}

\entry{Cornette}{}{}{À vérifier dans un dictionnaire lexicographique.}{}{Tous les cœurs se rallient à sa blanche cornette, Si le chrétien succombe à son charme insidieux,}

\entry{Cotart}{}{}{À vérifier dans un dictionnaire lexicographique.}{}{L'âme du bon feu maistre Jehan Cotart, se réincarnait chez ce vieux fêtard.}

\entry{Cotillon}{}{nom}{Jupe de femme, notamment dans un emploi ancien; danse populaire à l’origine.}{}{Quand il se fit tendre, elle lui dit : « J'présage, Qu'c'est pas dans les plis de mon cotillon}

\entry{Couillonnades}{}{}{À vérifier dans un dictionnaire lexicographique.}{}{Soit dit sans couillonnades avait le nom d'un ad-jectif démonstratif}

\entry{Coulpe}{}{nom}{Faute ou culpabilité, surtout dans « battre sa coulpe ».}{}{Le criminel saisi d'un zèle expiatoire, Qui bat sa coulpe bourrelé par le remords,}

\entry{Cousette}{}{}{À vérifier dans un dictionnaire lexicographique.}{}{Qui brode, divine cousette, des arcs-en-ciel à nos chaussettes ?}

\entry{Croquantes}{}{}{À vérifier dans un dictionnaire lexicographique.}{}{Toi qui m'as donné du feu quand, Les croquantes et les croquants}

\entry{Croquants}{}{}{À vérifier dans un dictionnaire lexicographique.}{}{Toi qui m'as donné du feu quand, Les croquantes et les croquants}

\entry{Croquignol}{}{adj./nom}{Ridicule, grotesque; personne ou situation comique par son aspect.}{}{Et donnèrent je vous l'assure, Un spectacle assez croquignol.}

\entry{Crottés}{}{}{À vérifier dans un dictionnaire lexicographique.}{}{Les sabots d'Hélène, étaient tout crottés,}

\entry{Crottin}{}{}{À vérifier dans un dictionnaire lexicographique.}{}{D'voir leurs héritiers marrons marcher dans le crottin}

\end{multicols}

%----------------------------------------------------------------------------------------
%	SECTION D
%----------------------------------------------------------------------------------------

\section*{D}

\begin{multicols}{3}

\extramarks{D}{D}
\entry{Danaïdes}{}{}{À vérifier dans un dictionnaire lexicographique.}{}{C'est le tonneau des Danaïdes changé en femme, Les joies charnelles m'emmerdent. \{2x\}}

\entry{Dandys}{}{}{À vérifier dans un dictionnaire lexicographique.}{}{Tous les Brummel, les dandys, les gandins, il les considérait avec dédain}

\entry{De concert}{}{}{En accord ou ensemble avec d’autres.}{}{}

\entry{De conserve}{}{}{Ensemble, en commun; souvent à propos de personnes ou de navires.}{}{}

\entry{De haute lutte}{}{}{Après une lutte difficile et acharnée.}{}{}

\entry{Délectable}{}{adj.}{Qui procure un grand plaisir, notamment au goût ou à l’esprit.}{}{La suite serait délectable, Malheureusement je ne peux,}

\entry{Dès potron-minet}{}{}{Très tôt le matin, à l’aube.}{}{}

\entry{Désempoissonnée}{}{}{À vérifier dans un dictionnaire lexicographique.}{}{Le crêpe à la queue sans doute, L'escorteront chagrinés, Laissant la rivière toute, Vide, désempoissonnée.}

\entry{Donzelle}{}{nom}{Jeune fille ou femme, terme ancien ou plaisant.}{}{Me sentant rempli de pitié, Pour la donzelle, J'lui enseignai, de son métier, Les p'tites ficelles}

\end{multicols}

%----------------------------------------------------------------------------------------
%	SECTION E
%----------------------------------------------------------------------------------------

\section*{E}

\begin{multicols}{3}

\extramarks{E}{E}
\entry{Émule}{}{}{À vérifier dans un dictionnaire lexicographique.}{}{Quand on ne le veut point, un émule de Socrate,}

\entry{En filigrane}{}{}{De façon discrète mais perceptible en arrière-plan.}{}{}

\entry{En un tournemain}{}{}{Très rapidement et avec facilité.}{}{N'y allant pas par quatre chemins, J'estourbis en un tournemain}

\entry{Engeance}{}{nom}{Race, espèce, descendance; souvent avec une nuance péjorative.}{}{Dès que la féminine engeance, sut que le singe était puceau}

\entry{Espérance}{}{}{À vérifier dans un dictionnaire lexicographique.}{}{J'ai l'espérance qu'elle viendra, Faire sa niche entre mes bras, entre mes bras.}

\entry{Être aux abois}{}{}{Être dans une situation désespérée, sans solution apparente.}{}{}

\entry{Exutoires}{}{}{À vérifier dans un dictionnaire lexicographique.}{}{Le bagne, l'échafaud entre autres exutoires, Et l'efficacité de la peine de mort,}

\end{multicols}

%----------------------------------------------------------------------------------------
%	SECTION F
%----------------------------------------------------------------------------------------

\section*{F}

\begin{multicols}{3}

\extramarks{F}{F}
\entry{Fâcheux}{}{adj./nom}{Qui importune, ennuie ou contrarie; personne importune.}{}{Mais une attention profonde, Prouve que c'est chez les fâcheux, Qu'il préfère choisir les victimes de ses petits jeux}

\entry{Faire bonne chère}{}{}{Bien manger; recevoir ou manger abondamment.}{}{}

\entry{Faire des gorges chaudes}{}{}{Se moquer abondamment d’une personne ou d’un événement.}{}{}

\entry{Faire fi de}{}{}{Ne pas tenir compte de; mépriser.}{}{}

\entry{Faire florès}{}{}{Connaître un grand succès, être en vogue.}{}{}

\entry{Faire la nique}{}{}{Se moquer de quelqu’un ou le défier par un geste de mépris.}{}{}

\entry{Faire la pige au temps}{}{}{Rivaliser avec le temps, tenter de lui échapper ou de le tromper.}{}{}

\entry{Faire long feu}{}{}{Échouer ou ne pas produire l’effet attendu; selon l’usage, durer plus longtemps que prévu.}{}{}

\entry{Faire route}{}{}{Se diriger vers un lieu, voyager dans une direction donnée.}{}{}

\entry{Faire un four}{}{}{Subir un échec, notamment pour un spectacle ou une entreprise.}{}{}

\entry{Fretin}{}{nom}{Petit poisson; au figuré, personnes ou choses de peu d’importance.}{}{Mais s'il pêche, c'est pour rire, Et l'on peut être certain, Que jamais sa poêle à frire, Vit le plus menu fretin.}

\entry{Froufrous}{}{}{À vérifier dans un dictionnaire lexicographique.}{}{La demeure que je préfère, C'est votre robe à froufrous}

\end{multicols}

%----------------------------------------------------------------------------------------
%	SECTION G
%----------------------------------------------------------------------------------------

\section*{G}

\begin{multicols}{3}

\extramarks{G}{G}
\entry{Gambades}{}{}{À vérifier dans un dictionnaire lexicographique.}{}{Lors ma mie, sans croire au danger, Faisons mille et une gambades,}

\entry{Gandins}{}{}{À vérifier dans un dictionnaire lexicographique.}{}{Tous les Brummel, les dandys, les gandins, il les considérait avec dédain}

\entry{Gargotière}{}{nom}{Femme qui tient une gargote.}{}{Madame ma gargotière, Comme je lui dois trop de sous}

\entry{Gastibelza}{}{}{À vérifier dans un dictionnaire lexicographique.}{}{GASTIBELZA "L'HOMME A LA CARABINE"}

\entry{Gendarmicides}{}{}{À vérifier dans un dictionnaire lexicographique.}{}{Des mégères gendarmicides, En criant: « Hip, hip, hip, hourra ! »}

\entry{Gentilhommière}{}{nom}{Petite demeure de campagne appartenant à un gentilhomme.}{}{Je changerais ma chaumière, Pour une gentilhommière}

\entry{Gentilhommières}{}{}{À vérifier dans un dictionnaire lexicographique.}{}{Où le septième ciel, ma plus chère ballade, Ma plus douce grimpette et plus tendre escalade, Sera trop haut pour moi. Il n'y aura pas de pleurs dans les gentilhommières, Ni de grincements de fesses dans les chaumières, Faut pas que je me leurre.}

\entry{Gibus}{}{}{À vérifier dans un dictionnaire lexicographique.}{}{J'ai des velléités d'arpenter les trottoir(e)s, Avec cette devise écrite à mon gibus :}

\entry{Gnons}{}{}{À vérifier dans un dictionnaire lexicographique.}{}{Jugeant enfin que leurs victimes, Avaient eu leur content de gnons}

\entry{Goton}{}{nom}{Femme ou fille de condition modeste; nom populaire et vieilli.}{}{Et plus belle d'une goton. Dieu, s'il existe, il exagère,}

\entry{Goualante}{}{nom}{Chanson populaire, souvent d’argot ou de rue.}{}{Poussant une autre goualante y'aurait à ma place un autre cabot}

\entry{Gougoutte}{}{}{À vérifier dans un dictionnaire lexicographique.}{}{Quand Margot dégrafait son corsage, pour donner la gougoutte à son chat}

\entry{Goupillon}{}{}{À vérifier dans un dictionnaire lexicographique.}{}{Et l'on courut à toutes jambes, Quérir un goupillon mais}

\entry{Gourgandines}{}{}{À vérifier dans un dictionnaire lexicographique.}{}{Mais s'il attrape une ondine, L'une de ces gourgandines, Femme mi-chair mi-poisson, Le polisson}

\entry{Grégeois}{}{}{À vérifier dans un dictionnaire lexicographique.}{}{Et l'Espagne flambait dans un grand feu grégeois. Je chantais, et j'étais pas le seul « Y a d'la joie ».}

\entry{Guilleret}{}{adj.}{Gai, enjoué, de bonne humeur.}{}{Laisse-moi tenir ton jupon, Courons guilleret, guillerette,}

\entry{Guillerette}{}{}{À vérifier dans un dictionnaire lexicographique.}{}{Laisse-moi tenir ton jupon, Courons guilleret, guillerette,}

\end{multicols}

%----------------------------------------------------------------------------------------
%	SECTION H
%----------------------------------------------------------------------------------------

\section*{H}

\begin{multicols}{3}

\extramarks{H}{H}
\entry{Halieutique}{}{adj.}{Relatif à la pêche, notamment à la pêche maritime.}{}{On dirait un fanatique, De la cause halieutique, Avec sa belle canne et, Son moulinet.}

\entry{Hottentots}{}{}{À vérifier dans un dictionnaire lexicographique.}{}{Callipyge à prétendre jouer les Vénus chez les Hottentots}

\end{multicols}

%----------------------------------------------------------------------------------------
%	SECTION I
%----------------------------------------------------------------------------------------

\section*{I}

\begin{multicols}{3}

\extramarks{I}{I}
\entry{Infamie}{}{}{À vérifier dans un dictionnaire lexicographique.}{}{Trouverais-je les noms, trouverais-je les mots, Pour noter d'infamie cet enfant de chameau}

\entry{Ingénu}{}{}{À vérifier dans un dictionnaire lexicographique.}{}{Faut s'lever de bon matin pour voir un ingénu, Qui n't'ait pas connue,}

\end{multicols}

%----------------------------------------------------------------------------------------
%	SECTION J
%----------------------------------------------------------------------------------------

\section*{J}

\begin{multicols}{3}

\extramarks{J}{J}
\entry{Jarnibleus}{}{}{À vérifier dans un dictionnaire lexicographique.}{}{Ainsi, parbleu, que les jarnibleus, Et les palsambleus}

\entry{Jarnicotons}{}{interj.}{Juron fantaisiste, variation populaire de jurons anciens.}{}{Sans oublier les jarnicotons, Les scrogneugneus et les bigres et les bougres}

\entry{Jobastre}{}{nom/adj.}{Personne extravagante, maladroite ou stupide; terme populaire.}{}{Si le sieur Z était un jobastre sans grade, Il laisserait en paix ses pauvres camarades.}

\entry{Jouer des fesses}{}{}{Bouger ou utiliser ses fesses de façon suggestive; dans le texte, se mouvoir de manière sexuelle.}{}{Car, dans l'art de faire le trottoir, Je le confesse, Le difficile est d'bien savoir, Jouer des fesses}

\entry{Jouvenceau}{}{nom}{Jeune homme, garçon encore très jeune.}{}{Depuis que je commence à faire de vieux os, Avide de conseils, souvent un jouvenceau}

\end{multicols}

%----------------------------------------------------------------------------------------
%	SECTION L
%----------------------------------------------------------------------------------------

\section*{L}

\begin{multicols}{3}

\extramarks{L}{L}
\entry{Lâcher la bride à son émoi}{}{}{Laisser libre cours à son émotion.}{}{}

\entry{Lacune}{}{}{À vérifier dans un dictionnaire lexicographique.}{}{C'est : « enculé ». Lacune comblée.}

\entry{Languedo}{}{}{À vérifier dans un dictionnaire lexicographique.}{}{Dans les années quarante où je débarquais de mon Languedo}

\entry{Lapidaire}{}{adj.}{Très concis, réduit à l’essentiel; relatif aussi à la taille ou au commerce des pierres.}{}{Ils gravent sur mon mur en style lapidaire : « Ici loge un vieux bouc qui n'a plus d'érections ! » Ils sont prématurés, tous ces cris de victoire, Ô vous qui me plantez la corne dans le dos,}

\entry{Licou}{}{nom}{Licol, dispositif de cuir ou de corde placé sur la tête d’un animal.}{}{Menait César, empereur d'Allemagne, Par le licou... -Le vent qui vient à travers la montagne Me rendra fou.}

\entry{Lorgner}{}{}{À vérifier dans un dictionnaire lexicographique.}{}{A l'heure du repas, mes rivaux détestables, Ont encor ce toupet de lorgner ma portion}

\entry{Lourdauds}{}{}{À vérifier dans un dictionnaire lexicographique.}{}{Une autre fourre avec rudesse, Le crâne d'un de ses lourdauds}

\entry{Lurette}{}{}{À vérifier dans un dictionnaire lexicographique.}{}{Voici belle lurette, en fut le vrai premier.}

\end{multicols}

%----------------------------------------------------------------------------------------
%	SECTION M
%----------------------------------------------------------------------------------------

\section*{M}

\begin{multicols}{3}

\extramarks{M}{M}
\entry{Macchabées}{}{}{À vérifier dans un dictionnaire lexicographique.}{}{Moi, j'bichais car je les adore, Sous la forme de macchabées}

\entry{Madone}{}{nom}{La Vierge Marie représentée dans l’art; par extension, représentation de la Vierge.}{}{J'lui ai dit: « De la Madone, Tu es le portrait ! » Le Bon Dieu me le pardonne, C'était un peu vrai...}

\entry{Maraud}{}{nom/adj.}{Personne grossière, malhonnête ou méprisable; parfois terme d’injure.}{}{Si, par hasard, Sur l'Pont des Arts, Tu croises le vent, le vent maraud, Prudent, prends garde à ton chapeau}

\entry{Maritornes}{}{nom}{Femmes grossières ou peu séduisantes; par référence au personnage de Maritorne dans Don Quichotte.}{}{Un truc, un moyen plausible, De fuir un peu son chez-soi, Où sévit la plus nuisible, Des maritornes qui soient.}

\entry{Matamore}{}{nom}{Fanfaron, vantard qui se donne des airs de héros.}{}{Mais il est général, va-t-en-guerre, matamore. Dès qu'il s'en mêle, on compte les morts.}

\entry{Matines}{}{}{À vérifier dans un dictionnaire lexicographique.}{}{Ding ding dong, les matines sonnent, En l'honneur de notre bonheur,}

\entry{Maugrabine}{}{nom}{Femme originaire du Maghreb ou, dans un emploi ancien, de l’Afrique du Nord.}{}{« Quelqu'un de vous a-t-il connu Sabine, Ma señora ? Sa mère était la vieille maugrabine D'Antequera,}

\entry{Mégères}{}{}{À vérifier dans un dictionnaire lexicographique.}{}{Des mégères gendarmicides, En criant: « Hip, hip, hip, hourra ! »}

\entry{Mener par le bout du cœur}{}{}{Exercer sur quelqu’un une forte emprise affective.}{}{}

\entry{Méritoire}{}{}{À vérifier dans un dictionnaire lexicographique.}{}{« Le meilleur d'entre nous et le plus méritoire », « Un saint homme, un cœur d'or, un bel et noble esprit »}

\entry{Métempsycose}{}{nom}{Passage de l’âme d’un corps dans un autre après la mort.}{}{Métamorphose inouïe, métempsycose infâme, Les joies charnelles me perdent,}

\entry{Mettre à feu et à sang}{}{}{Dévaster, ravager; causer de grands dégâts.}{}{}

\entry{Morbleu}{}{interj.}{Juronnement ancien, euphémisme de « mort de Dieu ».}{}{Souvenir de la cane, De Jeanne, Morbleu}

\entry{Morpionibus}{}{locution}{Formule latinisante et fantaisiste utilisée comme jeu verbal autour du mot « morpion ».}{}{Et quand il y alla de ses De Profondis, L'enfant de cœur répliqua pas morpionibus}

\end{multicols}

%----------------------------------------------------------------------------------------
%	SECTION N
%----------------------------------------------------------------------------------------

\section*{N}

\begin{multicols}{3}

\extramarks{N}{N}
\entry{Naturalibus}{}{locution}{Nu, sans vêtements, dans l’expression « in naturalibus ».}{}{Qui'est-ce qui veut m'laisser faire, in naturalibus, Un p'tit peu d'alpinisme sur son mont de Vénus}

\entry{Navrance}{}{nom}{État de grande tristesse, de désolation ou de découragement.}{}{On lirait quelle navrance mon blase inconnu dans un ex-voto}

\entry{Nectar}{}{}{À vérifier dans un dictionnaire lexicographique.}{}{Va boire à Passy, Le nectar d'ici, Te dépasse.}

\entry{Nonnettes}{}{}{À vérifier dans un dictionnaire lexicographique.}{}{De nonnettes, De cornettes, En sabbat,}

\end{multicols}

%----------------------------------------------------------------------------------------
%	SECTION O
%----------------------------------------------------------------------------------------

\section*{O}

\begin{multicols}{3}

\extramarks{O}{O}
\entry{Odalisques}{}{}{À vérifier dans un dictionnaire lexicographique.}{}{}

\entry{Onaniste}{}{}{À vérifier dans un dictionnaire lexicographique.}{}{La moitié des messieurs brûle d'être onaniste,}

\entry{Oraisons}{}{}{À vérifier dans un dictionnaire lexicographique.}{}{Sur les tombeaux les oraisons déclamatoires, Les : « C'était un bon fils, bon père, bon mari »,}

\entry{Oriflamme}{}{nom}{Bannière ou étendard, notamment historique; au figuré, symbole ou signe de ralliement.}{}{Il conquit un oriflamme teuton. Cet acte lui valut le grand cordon.}

\entry{Ostrogoth}{}{}{À vérifier dans un dictionnaire lexicographique.}{}{Au premier ostrogoth venu, Les croquants, ça tombe des nues}

\end{multicols}

%----------------------------------------------------------------------------------------
%	SECTION P
%----------------------------------------------------------------------------------------

\section*{P}

\begin{multicols}{3}

\extramarks{P}{P}
\entry{Palsambleus}{}{}{À vérifier dans un dictionnaire lexicographique.}{}{Ainsi, parbleu, que les jarnibleus, Et les palsambleus}

\entry{Pandores}{}{}{À vérifier dans un dictionnaire lexicographique.}{}{En voyant ces braves pandores, Être à deux doigts de succomber}

\entry{Pâquerette}{}{}{À vérifier dans un dictionnaire lexicographique.}{}{Indiscrète, Pâquerette, D'où vient-elle ?}

\entry{Pardi}{}{interj.}{Forme populaire de « par Dieu », employée pour renforcer une affirmation.}{}{Je n'perdais pas au change, pardi}

\entry{Pastourelle}{}{nom}{Petit poème ou chanson médiévale mettant en scène un berger ou une bergère; par extension, genre pastoral.}{}{Et là-dessus le Corydon, Le promis de la pastourelle,}

\entry{Patibulaire}{}{adj.}{Qui a une apparence inquiétante, suspecte ou peu rassurante.}{}{Je serais mort, jambes en l'air, Sur la veuve patibulaire,}

\entry{Pécore}{}{}{À vérifier dans un dictionnaire lexicographique.}{}{Fils de pécore et de minus, Ris par de la pauvre Vénus}

\entry{Peloteur}{}{nom}{Personne qui cherche à caresser ou à toucher quelqu’un de manière importune.}{}{Grand peloteur de fesses convaincu, passé maître en l'art de la main au cul,}

\entry{Pétaudière}{}{nom}{Lieu ou assemblée où règnent le désordre et la confusion.}{}{En mélangeant les genres, vous avez fait d'la terre, Ce qu'elle est : une pétaudière !}

\entry{Phallocrate}{}{}{À vérifier dans un dictionnaire lexicographique.}{}{Quand on veut les trousser, on est un phallocrate,}

\entry{Pimpant}{}{adj.}{Joli, bien mis, élégant ou plein de vivacité.}{}{Et tout pimpant, tout satisfait, Prit la clef du champ de navets.}

\entry{Pissotant}{}{}{À vérifier dans un dictionnaire lexicographique.}{}{Un soir d'intempérance, à son insu, il éteignit en pissotant dessus}

\entry{Populo}{}{}{À vérifier dans un dictionnaire lexicographique.}{}{A qui fera-t-on croire que le bon populo, quand il chante quand même, est un parfait salaud ?}

\entry{Pourceau}{}{}{À vérifier dans un dictionnaire lexicographique.}{}{Que le fameux cochon, le pourceau d'Epicure, Qui sommeillait en moi ne s'éveillera plus. Ils me croient interdit de séjour à Cythère, Et, par les nuits sans lune avec jubilation,}

\entry{Prendre la mort comme elle vient}{}{}{Accepter la mort sans chercher à la maîtriser.}{}{Et jamais je ne parviens, A prendre la mort comme elle vient,}

\entry{Preste}{}{adj.}{Rapide et agile dans ses mouvements.}{}{Le facteur d'ordinaire si preste pour voir ça ne distribuait plus les lettres que personne au reste n'aurait lues}

\entry{Priape}{}{}{À vérifier dans un dictionnaire lexicographique.}{}{Peu de chances qu'on voie mes belles odalisques, Déposer en grand deuil au pied de l'obélisque, Quelques gerbes de fleurs. Tout au plus gentiment diront-elles : « Peuchère, Le vieux Priape est mort », et, la cuisse légère, Le regard alangui,}

\entry{Probes}{}{}{À vérifier dans un dictionnaire lexicographique.}{}{Les jean-foutre et les gens probes, Médisent du vent furibond, Qui rebrousse les bois, détrousse les toits, retrousse les robes}

\end{multicols}

%----------------------------------------------------------------------------------------
%	SECTION Q
%----------------------------------------------------------------------------------------

\section*{Q}

\begin{multicols}{3}

\extramarks{Q}{Q}
\entry{Quadrumane}{}{nom/adj.}{Qui a quatre membres préhensiles; ancien terme appliqué notamment aux singes.}{}{Voyant que toutes se dérobe, Le quadrumane accéléra,}

\entry{Quenouille}{}{nom}{Bâton servant à filer les fibres textiles; par extension, objet ayant une forme allongée.}{}{Comme il atteignait l'orée du village, Filant sa quenouille, il vit Cendrillon}

\end{multicols}

%----------------------------------------------------------------------------------------
%	SECTION R
%----------------------------------------------------------------------------------------

\section*{R}

\begin{multicols}{3}

\extramarks{R}{R}
\entry{Réquisitoire}{}{}{À vérifier dans un dictionnaire lexicographique.}{}{Fatigué de souffrir leur long réquisitoire}

\entry{Retourner à ses oignons}{}{}{Revenir à ses propres affaires; se mêler de ce qui nous concerne.}{}{}

\entry{Ritournelles}{}{}{À vérifier dans un dictionnaire lexicographique.}{}{Le feu de la ville éternelle est éternel. Si Dieu veut l'incendie, il veut les ritournelles.}

\entry{Rogaton}{}{nom}{Objet sans valeur, petit reste ou écrit insignifiant; anciennement, reste de nourriture.}{}{Sans laisser même un rogaton, Tout il croque, tout il digère.}

\end{multicols}

%----------------------------------------------------------------------------------------
%	SECTION S
%----------------------------------------------------------------------------------------

\section*{S}

\begin{multicols}{3}

\extramarks{S}{S}
\entry{Sacrebleu}{}{}{À vérifier dans un dictionnaire lexicographique.}{}{Bande à part, sacrebleu ! c'est ma règle et j'y tiens. Dans les noms des partants on n'verra pas le mien.}

\entry{Sacristain}{}{nom}{Personne chargée de la garde et de l’entretien d’une église et de ses objets liturgiques.}{}{On n'tortille pas son popotin, D'la même manière, Pour un droguiste, un sacristain, Un fonctionnaire}

\entry{Sans ambages}{}{}{Directement, sans détour.}{}{Sans ambages, Dans les pages, Du missel,}

\entry{Sans coup férir}{}{}{Sans rencontrer de résistance.}{}{}

\entry{Sans vergogne}{}{}{Sans honte ni gêne.}{}{T'as couru sans vergogne, et pour une escalope, Te jeter dans le lit du boucher.}

\entry{Sbires}{}{}{À vérifier dans un dictionnaire lexicographique.}{}{Deux sbires sont venus, avec leur noirs manteaux, entre la rue Didot et la rue de Vanves}

\entry{Séguedille}{}{nom}{Danse et chant espagnols d’origine populaire, à rythme vif.}{}{Un pied pour la séguedille, un pied pour, La gaie tarentelle.}

\entry{Sinécure}{}{}{À vérifier dans un dictionnaire lexicographique.}{}{Parc' que depuis tant d'années, C'était pas une sinécure}

\entry{Subito}{}{}{À vérifier dans un dictionnaire lexicographique.}{}{La souris grise se fâche, et subito presto, entre la rue Didot et la rue de Vanves}

\entry{Suffragettes}{}{}{À vérifier dans un dictionnaire lexicographique.}{}{A force d'être en butte au tir des suffragettes}

\entry{Supplique}{}{}{À vérifier dans un dictionnaire lexicographique.}{}{05 - Supplique pour être enterré sur la plage de Sète\_sans\_accords}

\end{multicols}

%----------------------------------------------------------------------------------------
%	SECTION T
%----------------------------------------------------------------------------------------

\section*{T}

\begin{multicols}{3}

\extramarks{T}{T}
\entry{Tarabuste}{}{}{À vérifier dans un dictionnaire lexicographique.}{}{C'est très inconfortable et ça vous tarabuste,}

\entry{Tendre la patte}{}{}{Demander ou recevoir de l’argent ou une aide; selon le contexte, solliciter une faveur.}{}{}

\entry{Tenir la dragée haute}{}{}{Résister avec succès à quelqu’un, lui donner du fil à retordre.}{}{}

\entry{Têtes chenues}{}{}{Têtes blanchies par l’âge; image traditionnelle de la vieillesse.}{}{Quand ils sont dev'nus des têtes chenues, des grisons}

\entry{Teuton}{}{adj./nom}{Allemand, en emploi familier ou péjoratif, par référence aux anciens Teutons.}{}{Il conquit un oriflamme teuton. Cet acte lui valut le grand cordon.}

\entry{Tout de go}{}{}{Directement, sans détour ni préparation.}{}{Le chat la prenant pour sa mère se mit à téter tout de go, émue Margot le laissa faire brave Margot}

\entry{Tramontane}{}{nom}{Vent froid et sec venant du nord-ouest dans le Midi de la France; au figuré, idée de vent ou de perte de repères.}{}{J'ai perdu la tramontane En trouvant Margot, Princesse vêtue de laine, Déesse en sabots...}

\entry{Transcendant}{}{}{À vérifier dans un dictionnaire lexicographique.}{}{Taisons-le, j'aurais bonne mine. Il me paraît plus transcendant}

\entry{Transi}{}{}{À vérifier dans un dictionnaire lexicographique.}{}{Malheureux, poussiéreux, transi, Chante : « Ah ! ce qu'on s'emmerde ici » !}

\end{multicols}

%----------------------------------------------------------------------------------------
%	SECTION V
%----------------------------------------------------------------------------------------

\section*{V}

\begin{multicols}{3}

\extramarks{V}{V}
\entry{Valseur}{}{}{À vérifier dans un dictionnaire lexicographique.}{}{La louche au valseur; pas de ça Lisette ! La légion d'honneur ça pardonne pas.}

\entry{Vauvert}{}{nom propre / locution}{Nom historique du Val-de-Green; présent dans l’expression « au diable Vauvert ».}{}{Emportent les trépassés jusqu'au diable vauvert}

\entry{Ventrebleu}{}{}{À vérifier dans un dictionnaire lexicographique.}{}{Avec qui, ventrebleu, faut-il que je couche, Pour faire parler un peu la déesse aux cent bouches}

\entry{Vêpres}{}{nom}{Office religieux célébré en fin d’après-midi ou le soir.}{}{La mignonne allait aux vêpres, Se mettre à genoux, Alors j'ai mordu ses lèvres, Pour savoir leur goût...}

\entry{Vertudieux}{}{adj.}{Qui affecte ou manifeste une vertu morale ostentatoire.}{}{Tous les Bon Dieu, Tous les vertudieux, Tonnerre de Brest et saperlipopette}

\entry{Viatique}{}{nom}{Sacrement ou secours spirituel donné à un mourant; au figuré, ressource emportée pour un voyage.}{}{Pour offrir à l'ancêtre, en signe d'affection, En guise de viatique, une ultime audition (bis)}

\entry{Vicaire}{}{}{À vérifier dans un dictionnaire lexicographique.}{}{Avant même que le vicaire, Ait pu lâcher un cri}

\entry{Vocables}{}{nom}{Mots, termes ou expressions considérés sous leur aspect lexical.}{}{Brillait par son absence un des pires vocables}

\entry{Voir des deux yeux}{}{}{Constater directement, par opposition à croire sur parole.}{}{}

\end{multicols}

